\section{E3: screening additional checkpoints}
\label{app:e3}
E3 tests whether the replay model that agreed with E2 can screen additional
checkpoints before deployment. From the existing training archive, we replay
90 previously undeployed candidates under C and C30 on both locked windows
(360 CPU scenarios), without retraining. C accepts 67 candidates and rejects
23; C30 accepts 89. Acceptance requires P1+ on both windows. These are replay
judgments, not measured success rates.

\paragraph{Selection and execution.}
The selection rule uses C alone. Accepted candidates are ordered by mean
window cost; rejected candidates by total offline and then online deadline
misses, with checkpoint identity breaking ties. Within each group, identical
complete target trajectories are represented once, and the first three
representatives are selected. The resulting six C trajectories differ from
the previously deployed policies. Table~\ref{tab:e3identities} maps the new
labels to archived checkpoints; the selected set spans two training seeds.
C30 accepts all six but does not determine their selection.

\begin{table}[t]
\centering
\small
\caption{E3 checkpoint identities. Names encode training model, seed, and
checkpoint index. C rejects N4--N6 because it predicts respectively 3, 6,
and 7 offline misses in steady; all six pass recovery in replay.}
\label{tab:e3identities}
\setlength{\tabcolsep}{3pt}
\par\vspace{4pt}
\begin{tabular*}{\linewidth}{@{\extracolsep{\fill}}lll@{}}
\toprule
Arm & Archived checkpoint & C decision\\
\midrule
N1 & \texttt{C30\_s105\_e0065} & accept\\
N2 & \texttt{C30\_s105\_e0035} & accept\\
N3 & \texttt{C30\_s102\_e0005} & accept\\
N4 & \texttt{C30\_s102\_e0080} & reject\\
N5 & \texttt{C\_s105\_e0080} & reject\\
N6 & \texttt{C\_s105\_e0070} & reject\\
\bottomrule
\end{tabular*}
\end{table}

Each checkpoint executes twice per window, with one all2 reference per
window: 24 policy runs and two references. A checkpoint is physically
feasible only if all four executions satisfy P1+. Recovery and steady
execute on separate matched dual-A100 hosts with the same model weights,
vLLM 0.9.0 software stack, and accounting specifications. The recorded
execution amendment changes work allocation, retaining the six checkpoints,
inputs, predictions, and within-window run order. Window and host are
therefore confounded; differences between the windows cannot be attributed
solely to workload. E3 remains separate from the E2 runtime group.

The launch records bind the frozen scientific files and CPU predictions;
actions and replay outcomes match the CPU references. The external Git
record was created after physical execution began, so this is a fixed-cohort
confirmation, not an externally preregistered study. All 26 executions are
technically valid, with no retries or replacement candidates.

\paragraph{Results and interpretation.}
All 26 runs satisfy P1+. Thus C correctly accepts three checkpoints but
falsely rejects the other three; C30 accepts all six correctly on this
cohort. There are no physically infeasible checkpoints, and C30 has no
rejected candidates, so neither failure-detection performance nor C30's
rejection accuracy can be estimated. The selected cohort and repeated runs
also do not provide a population screening-accuracy estimate.

\begin{table*}[t]
\centering
\small
\caption{E3 window costs in scenario USD, with both physical repeats shown.
R and S denote recovery and steady. Predictions include any modeled deadline
penalty; no physical run incurs one. Each all2 reference satisfies P1+ with
cost 4.2000. Costs cover the 2100\,s observation window; setup and subsequent
cleanup are recorded separately.}
\label{tab:e3costs}
\setlength{\tabcolsep}{4pt}
\par\vspace{4pt}
\begin{tabular*}{0.94\linewidth}{@{\extracolsep{\fill}}llrrrrr@{}}
\toprule
Arm & Window & C prediction & C30 prediction & Physical 1 & Physical 2 & P1+\\
\midrule
N1 & R & 1.9351 & 1.9351 & 1.9269 & 1.9270 & 2/2\\
N1 & S & 1.9767 & 1.9668 & 1.9575 & 1.9579 & 2/2\\
N2 & R & 1.9102 & 1.8709 & 1.8698 & 1.8445 & 2/2\\
N2 & S & 2.0171 & 1.9517 & 1.9236 & 1.9275 & 2/2\\
N3 & R & 1.9720 & 1.9468 & 1.9315 & 1.9322 & 2/2\\
N3 & S & 2.0323 & 1.9760 & 1.9617 & 1.9618 & 2/2\\
N4 & R & 1.9650 & 1.9621 & 1.9523 & 1.9524 & 2/2\\
N4 & S & 2.0689 & 2.0180 & 2.0171 & 2.0172 & 2/2\\
N5 & R & 2.6433 & 3.6908 & 3.6664 & 3.6667 & 2/2\\
N5 & S & 3.0121 & 3.6548 & 3.6242 & 3.6243 & 2/2\\
N6 & R & 2.8257 & 3.7802 & 3.7660 & 3.7580 & 2/2\\
N6 & S & 3.0603 & 3.9023 & 3.8758 & 3.8769 & 2/2\\
\bottomrule

\end{tabular*}
\end{table*}

Across the 24 policy runs, feasibility agreement is 18/24 for C and
24/24 for C30: C predicts failure only in the steady cells of N4--N6.
Cost MAE is 0.3112 for C and 0.0163 for C30
(Table~\ref{tab:e3costs}); the two all2 references are excluded from this
comparison. These are scenario costs, including the declared synthetic-API
route, rather than measured provider charges. The outcome reverses the
model ordering observed in E2 on a different cohort and execution context;
it does not establish that either fixed model is generally preferable.

In this execution context, the 76 completed in-window launch-to-ready
intervals have median 34.78\,s
(range 32.52--93.01\,s), closer to C30's 30\,s assumption than C's historical
69.219\,s value. Two additional late launches are terminated at the
observation cutoff before readiness and are excluded from this duration
summary as right-censored observations; their window occupancy is charged.
These observations do not isolate startup time as the cause of the gate
errors. E3 shows why a screening decision itself requires physical
validation: a model successful on an earlier cohort can reject feasible
policies in a later deployment context.
